% ==============================================================================
% 00_conclusion.tex — UE SIS589
% ==============================================================================
\chapter*{Conclusion Générale}
\addcontentsline{toc}{chapter}{Conclusion Générale}
\label{chap:conclusion}

% ==============================================================================
\section*{Ce que vous avez construit}
% ==============================================================================

Au terme de ce cours, LiZbiZ-SIS n'est plus l'organisation aveugle du
chapitre d'introduction. Depuis le niveau~1 de maturité SOC — aucun
SIEM, aucune visibilité, un MTTD supérieur à trois semaines — vous avez
construit, brique par brique, une capacité de supervision opérationnelle~:
des journaux collectés en temps réel, des règles de corrélation mappées
sur MITRE ATT\&CK, une architecture de détection en profondeur (NIDS
périmétrique, HIDS sur chaque hôte, EDR en test), des playbooks rédigés,
une procédure d'acquisition forensique conforme à ISO~27037, et un
processus de réponse aux incidents ancré dans le droit camerounais.

Ce n'est pas un état final. C'est le niveau~3 d'un parcours qui continue.

% ==============================================================================
\section*{Les trois convictions de ce cours}
% ==============================================================================

\textbf{Première conviction~:} la sécurité opérationnelle est une
discipline de rigueur documentaire autant que de compétence technique.
Un dump mémoire sans PV de collecte est inutilisable en justice.
Une alerte sans ticket traçable ne sert à rien. Un incident sans RETEX
se répétera. La forme n'est pas du formalisme~: c'est la condition
de l'efficacité.

\textbf{Deuxième conviction~:} le contexte africain et camerounais
n'est pas un handicap à compenser, c'est un environnement à maîtriser.
La Loi~2010/012, la Convention de Malabo, les spécificités des opérateurs
mobile money, la réalité des infrastructures hybrides cloud et on-premise~:
tout cela définit le terrain réel sur lequel vous opérerez. Les outils
sont globaux~; le contexte d'application est local.

\textbf{Troisième conviction~:} les outils open-source (Wazuh, ELK,
Volatility, The Sleuth Kit, Autopsy, MISP) permettent de construire
des capacités défensives de niveau professionnel avec des budgets
d'organisation de taille intermédiaire. La contrainte budgétaire ne
justifie pas l'absence de supervision. Elle impose la créativité
et la priorisation.

% ==============================================================================
\section*{La suite}
% ==============================================================================

SIS589 vous a donné les fondamentaux. La maîtrise vient avec la pratique
continue~:

\begin{itemize}
  \item \textbf{Certifications de référence~:} CompTIA Security+,
        GIAC GCFE / GCIH / GREM (SANS Institute), eJPT / eCDFP (eLearnSecurity),
        certifications ISACA (CISA, CISM, CRISC).
  \item \textbf{Plateformes de pratique~:} TryHackMe (Blue Team Labs),
        CyberDefenders, LetsDefend, Boss of the SOC (Splunk).
  \item \textbf{Ressources~:} SANS Reading Room, DFIR.training,
        MITRE ATT\&CK navigator, Elastic Security Labs,
        Mandiant M-Trends annual report.
  \item \textbf{Communauté~:} CERT-MU (Cameroun), AfricaCERT,
        FIRST (Forum of Incident Response and Security Teams),
        groupes ISACA Afrique centrale.
\end{itemize}

\begin{boxresume}
\textbf{L'essentiel~:} les attaquants s'adaptent en permanence.
La veille, la formation continue et la pratique sur cas réels ne
sont pas des options~; elles sont la définition même du métier
d'analyste SOC ou d'investigateur numérique. Un SOC qui ne s'améliore
pas régresse.
\end{boxresume}

\vspace{3ex}
\begin{center}
  {\color{CouleurAccent}\rule{0.55\linewidth}{1pt}}\\[1ex]
  {\small\itshape\color{CouleurPrimaire}
    Rédigé avec la conviction que la maîtrise du numérique\\
    est une responsabilité collective autant qu'individuelle.}\\[0.7ex]
  {\small\bfseries\color{CouleurPrimaire}\auteurcours}\\
  {\footnotesize\color{CouleurBordInfo}\itshape
    \institutionprimaire{} --- \anneescolaire}
\end{center}
